\documentclass[11pt]{article}

\usepackage[margin=1in]{geometry}
\usepackage{amsmath,amssymb}
\usepackage{array}
\usepackage{booktabs}
\usepackage{enumitem}
\usepackage{xcolor}
\usepackage{hyperref}
\usepackage{setspace}

\setstretch{1.10}
\setlength{\emergencystretch}{2em}

\hypersetup{
    colorlinks=true,
    linkcolor=blue!50!black,
    urlcolor=blue!50!black,
    citecolor=black
}

% --------------------
% Course information
% --------------------
\newcommand{\course}{DSA4213: Natural Language Processing for Data Science}
\newcommand{\semester}{Semester 1, AY2026/2027}

\begin{document}

\begin{center}
{\LARGE\bfseries Group Project Proposal}\\[0.4em]
{\large Project Title}\\[0.9em]
Group ID: 18 \underline{\hspace{3cm}}\\[0.9em]
\course\\
\semester
\end{center}

% ====================
% Administrative information
% ====================
\section*{Administrative Information}

\begin{tabular}{@{}>{\bfseries\raggedright\arraybackslash}p{0.25\linewidth}
                    >{\raggedright\arraybackslash}p{0.70\linewidth}@{}}
Group members & Name (Student ID) for each member. \\[0.7em]

Chua Yong Sheng Joel (A0282307H)
Ng Jia Hao Sherwin (A0257932W)
Robin Ghosh (A0271671A)
Chong Yu Xuan (A0272687M)


Project direction & One short phrase describing the project area. \\[0.7em]
Project direction: Agentic AI + Judge-Aware Evaluation 
Project Area: Many-to-One Prompt Injection 

Repository & Link, if available: \url{https://github.com/LeuRNburgJoe11/DSA4213_Group18} \\
\end{tabular}

% ====================
% Executive summary
% ====================
\section*{Project Summary}

Summarize the agent-related question, the proposed system or study, the main
evaluation plan, and what the project aims to learn or demonstrate.

Our project aims to stimulate 3 agents prompt injecting into another agent and find out which agent is effective in detecting and resisting prompt injection.
We also aim to 

% ====================
% 1. Problem and motivation
% ====================
\section*{1. Problem, Research Question, and Motivation}

State one clear question, objective, or hypothesis. Explain the practical or
scientific motivation, the intended users or setting, and why an \emph{agent-based}
approach is relevant. Define what success would mean for this project.

% ====================
% 2. Related work and contribution
% ====================
\section*{2. Related Work and Proposed Contribution}

Identify the most relevant prior systems, datasets, benchmarks, papers, or tools.
Explain what the group will contribute beyond assembling existing components.
Cite the sources on which the project design depends.

% ====================
% 3. Agent setup
% ====================
\section*{3. Agent and Task Setup}

Describe the proposed setup precisely. For an empirical study, describe the agent
system and the study design or conditions considered. The list below is a guide,
not a checklist: you do not need to address every item. Include only the elements
that are relevant to your project:

\begin{itemize}[leftmargin=*,itemsep=0.25em]
    \item task and operating environment;
    \item inputs or observations and outputs or actions;
    \item model(s), prompts, tools, APIs, retrieval sources, memory, or other agents;
    \item users, simulators, datasets, benchmarks, or interaction protocols; and
    \item stopping conditions, constraints, permissions, and failure recovery.
\end{itemize}

% ====================
% 4. System or study design
% ====================
\section*{4. Proposed System or Study Design}

Describe what the group will implement or investigate and the sequence of the main
components. Identify the part that is technically meaningful and owned by the group.
If useful, include a system diagram, pseudocode, or an example agent trajectory.

\subsection*{4.1 Comparisons and Ablations (if applicable)}

Describe any baselines, controlled comparisons, or ablations that would help answer
the project question.

% ====================
% 5. Evaluation
% ====================
\section*{5. Preliminary Evaluation Plan}

\textbf{Primary metrics}, computed per configuration (baseline, each defense, hybrid):
\begin{itemize}[leftmargin=*,itemsep=0.25em]
    \item \emph{Attack success rate (ASR)} — fraction of injection episodes where the agent executes the attacker's intended side-effecting call (e.g., sends to the wrong recipient).
    \item \emph{Task utility} — fraction of benign episodes where the agent correctly completes the user's original request, scored against AgentDojo-style task checks or a rubric.
    \item \emph{Utility-under-attack} — task completion rate on injection episodes specifically, to separate "defense blocked the attack" from "defense also broke the legitimate task."
\end{itemize}

\textbf{Secondary/diagnostic metrics:}
\begin{itemize}[leftmargin=*,itemsep=0.25em]
    \item Tool-call budget used per episode (proxy for latency/overhead added by each defense).
    \item False-positive rate of the confirmation-gate classifier on benign side-effecting calls.
    \item Breakdown of ASR by injection type (visible / obfuscated / multi-step delayed-trigger) and position (early / late), tying results back to the section 4.1 ablations.
\end{itemize}

\textbf{Protocol:}
\begin{itemize}[leftmargin=*,itemsep=0.25em]
    \item Fixed set of benign + injection task instances, run identically across all four configurations for a controlled comparison.
    \item Multiple runs per episode (e.g., $n=3$–$5$) to account for LLM sampling variance; report mean $\pm$ spread rather than a single run.
    \item Automated pass/fail checks on final tool-call logs where possible; manual review for ambiguous trajectories.
\end{itemize}

\textbf{Qualitative analysis:} 2–3 representative successful defenses (attack blocked, task still completed) and 2–3 representative failures (attack succeeded, or defense broke utility) included as annotated trajectories in the report to illustrate why a defense worked or failed, not just that it did.

\textbf{Presentation:} results as a summary table (configuration $\times$ ASR $\times$ utility) plus one plot showing the security-utility trade-off frontier across configurations, supplemented by the trajectory case studies.
% ====================
% 6. Risks and responsible practice
% ====================
\section*{6. Risks, Limitations, and Responsible Practice}

Identify the main risks or limitations relevant to the project. Where appropriate,
describe safeguards or a fallback. Do not include credentials, private data, or
material that the group does not have permission to share.

\begin{itemize}[leftmargin=*,itemsep=0.25em]
    \item \textbf{Scope limitation:} findings are specific to the simulated environment and chosen base LLM(s); ASR and utility numbers may not transfer to other tool sets, models, or real deployments — this will be stated explicitly as a limitation, not generalized.
    \item \textbf{Injection realism:} synthetic injection prompts may not reflect the sophistication of real-world attacks; the group will note this as a threat-model boundary rather than claim comprehensive coverage.
    \item \textbf{Model/API variance:} LLM outputs are non-deterministic and API models may change over the project timeline; mitigated by pinning model versions where possible and running multiple trials per episode.
    \item \textbf{Dual-use consideration:} the project involves designing working prompt-injection attacks. These will be developed and tested strictly within the group's own sandboxed simulation, never against any real third-party service, and will be presented alongside their corresponding defenses rather than as standalone exploits.
    \item \textbf{No real data:} all task content (emails, calendar entries, etc.) will be synthetic — no real personal data, credentials, or private material will be used or included in the report or repository.
    \item \textbf{Fallback plan:} if API costs or time constraints limit the number of defenses/ablations that can be run, the group will prioritize the baseline vs. two defenses and the single-turn vs. delayed-trigger comparison, dropping remaining ablations first.
\end{itemize}




\bibliographystyle{plain}
\bibitem{debenedetti2024agentdojo}
E.~Debenedetti, J.~Zhang, M.~Balunovi\'{c}, L.~Beurer-Kellner, M.~Fischer, and
F.~Tram\`{e}r.
\newblock AgentDojo: A dynamic environment to evaluate attacks and defenses
for LLM agents.
\newblock \emph{NeurIPS Datasets and Benchmarks Track}, 2024.

\bibitem{greshake2023youve}
K.~Greshake, S.~Abdelnabi, S.~Mishra, C.~Endres, T.~Holz, and M.~Fritz.
\newblock Not what you've signed up for: Compromising real-world LLM-integrated
applications with indirect prompt injection.
\newblock \emph{Proceedings of the 16th ACM Workshop on Artificial
Intelligence and Security (AISec)}, 2023.

\bibitem{chen2024struq}
S.~Chen, J.~Piet, C.~Sitawarin, and D.~Wagner.
\newblock StruQ: Defending against prompt injection with structured queries.
\newblock \emph{arXiv preprint arXiv:2402.06363}, 2024.
% \bibliography{references}

\end{document}
